% GENERATED FILE. Edit canonical lessons, then run npm run generate:books.
% canonical-source-hash: fnv1a64:25e7b1cd6dce4e33
% canonical-lessons: MR-C90-bhint, MR-C90-chhappar, MR-C90-palang, MR-C90-kapat, MR-C90-jamin

\chapter{Wall, Roof, Bed, Cupboard, Ground}
\label{ch:mr-wall-roof-bed-cupboard-ground}
\hlchaptermodality{90}

\begin{chapteropening}
\textbf{By the end of this chapter:} \emph{I can say a wall, a roof, a bed, a cupboard and the ground, the floor.}

Complete the last lesson of chapter 90: I can say a wall, a roof, a bed, a cupboard and the ground, the floor.
\end{chapteropening}

\section[\texorpdfstring{\mr{भिंत} (bhint)}{भिंत (bhint)}]{\mr{भिंत} (bhint) — a wall}

\label{lesson:MR-C90-bhint}



\begin{quote}
Before the new one: say the Marathi for a chair, then the Marathi for a table.
\end{quote}

\subsection*{You'll want to know: \mr{भिंत}}
\textbf{\mr{भिंत}} — \emph{bhint} — \textquotedblleft{}a wall\textquotedblright{}.

\begin{cousinweb}
\textbf{\mr{भिंत}} (\emph{bhint}) is a tadbhava, a word inherited from Sanskrit and reshaped by sound change. It comes from \emph{bhitti}, \textquotedblleft{}mat of split reeds; wall\textquotedblright{} (CDIAL 9494). It is built on \emph{bhid}, \textquotedblleft{}to split\textquotedblright{}.
\end{cousinweb}

\begin{grammarlens}[title={What you've built}]
One more word for this chapter.
\end{grammarlens}

\subsection*{Your turn}
\begin{itemize}
\raggedright
  \item \emph{Say it:} \emph{bhint}
  \item \emph{Say it:} \emph{bhint}, once more
  \item \emph{Say it:} say \emph{bhint}
  \item \emph{Recall:} say the Marathi for a chair, then the Marathi for a table, then say \emph{bhint} again
\end{itemize}

\begin{tcolorbox}[breakable,colback=teal!4,colframe=teal!35!black,title={Before you move on}]
What does \mr{भिंत} mean? (\textquotedblleft{}A wall\textquotedblright{}.) Say it once more.
\end{tcolorbox}
\section[\texorpdfstring{\mr{छप्पर} (chhappar)}{छप्पर (chhappar)}]{\mr{छप्पर} (chhappar) — a roof}

\label{lesson:MR-C90-chhappar}



\begin{quote}
Before the new one: say the Marathi for a table, then the Marathi for a wall.
\end{quote}

\subsection*{You'll want to know: \mr{छप्पर}}
\textbf{\mr{छप्पर}} — \emph{chhappar} — \textquotedblleft{}a roof\textquotedblright{}.

\begin{cousinweb}
The origin of \textbf{\mr{छप्पर}} (\emph{chhappar}) is not settled. CDIAL has no entry for it; it may belong with a reconstructed \emph{*chapp-}, \textquotedblleft{}press, cover, hide\textquotedblright{} (CDIAL 4994), but no source says so.
\end{cousinweb}

\begin{grammarlens}[title={What you've built}]
One more word for this chapter.
\end{grammarlens}

\subsection*{Your turn}
\begin{itemize}
\raggedright
  \item \emph{Say it:} \emph{chhappar}
  \item \emph{Say it:} \emph{chhappar}, once more
  \item \emph{Say it:} say \emph{chhappar}
  \item \emph{Recall:} say the Marathi for a table, then the Marathi for a wall, then say \emph{chhappar} again
\end{itemize}

\begin{tcolorbox}[breakable,colback=teal!4,colframe=teal!35!black,title={Before you move on}]
What does \mr{छप्पर} mean? (\textquotedblleft{}A roof\textquotedblright{}.) Say it once more.
\end{tcolorbox}
\section[\texorpdfstring{\mr{पलंग} (palaṅg)}{पलंग (palaṅg)}]{\mr{पलंग} (palaṅg) — a bed}

\label{lesson:MR-C90-palang}



\begin{quote}
Before the new one: say the Marathi for a wall, then the Marathi for a roof.
\end{quote}

\subsection*{You'll want to know: \mr{पलंग}}
\textbf{\mr{पलंग}} — \emph{palaṅg} — \textquotedblleft{}a bed\textquotedblright{}.

\begin{cousinweb}
\textbf{\mr{पलंग}} (\emph{palaṅg}) goes back to Sanskrit \emph{palyaṅka}, \textquotedblleft{}bed\textquotedblright{} (CDIAL 7964). CDIAL says the north-western form with \emph{-ṅg} was borrowed widely as the word for \textquotedblleft{}bed\textquotedblright{}; Marathi's own inherited forms are \emph{pālak}, \textquotedblleft{}cradle\textquotedblright{}, and \emph{pālkhī}, \textquotedblleft{}litter\textquotedblright{}.
\end{cousinweb}

\begin{grammarlens}[title={What you've built}]
One more word for this chapter.
\end{grammarlens}

\subsection*{Your turn}
\begin{itemize}
\raggedright
  \item \emph{Say it:} \emph{palaṅg}
  \item \emph{Say it:} \emph{palaṅg}, once more
  \item \emph{Say it:} say \emph{palaṅg}
  \item \emph{Recall:} say the Marathi for a wall, then the Marathi for a roof, then say \emph{palaṅg} again
\end{itemize}

\begin{tcolorbox}[breakable,colback=teal!4,colframe=teal!35!black,title={Before you move on}]
What does \mr{पलंग} mean? (\textquotedblleft{}A bed\textquotedblright{}.) Say it once more.
\end{tcolorbox}
\section[\texorpdfstring{\mr{कपाट} (kapāṭ)}{कपाट (kapāṭ)}]{\mr{कपाट} (kapāṭ) — a cupboard}

\label{lesson:MR-C90-kapat}



\begin{quote}
Before the new one: say the Marathi for a roof, then the Marathi for a bed.
\end{quote}

\subsection*{You'll want to know: \mr{कपाट}}
\textbf{\mr{कपाट}} — \emph{kapāṭ} — \textquotedblleft{}a cupboard\textquotedblright{}.

\begin{cousinweb}
\textbf{\mr{कपाट}} (\emph{kapāṭ}) is Sanskrit \emph{kapāṭa}, \textquotedblleft{}leaf of a door\textquotedblright{}, a tatsama; CDIAL lists it as a variant of \emph{kavāṭa}, whose inherited Marathi form is \emph{kavāḍ}, \textquotedblleft{}door\textquotedblright{} (CDIAL 2963). CDIAL notes the view that the word is perhaps from Dravidian. In Marathi it names a cupboard.
\end{cousinweb}

\begin{grammarlens}[title={What you've built}]
One more word for this chapter.
\end{grammarlens}

\subsection*{Your turn}
\begin{itemize}
\raggedright
  \item \emph{Say it:} \emph{kapāṭ}
  \item \emph{Say it:} \emph{kapāṭ}, once more
  \item \emph{Say it:} say \emph{kapāṭ}
  \item \emph{Recall:} say the Marathi for a roof, then the Marathi for a bed, then say \emph{kapāṭ} again
\end{itemize}

\begin{tcolorbox}[breakable,colback=teal!4,colframe=teal!35!black,title={Before you move on}]
What does \mr{कपाट} mean? (\textquotedblleft{}A cupboard\textquotedblright{}.) Say it once more.
\end{tcolorbox}
\section[\texorpdfstring{\mr{जमीन} (jamīn)}{जमीन (jamīn)}]{\mr{जमीन} (jamīn) — the ground, the floor}

\label{lesson:MR-C90-jamin}



\begin{quote}
Before the new one: say the Marathi for a bed, then the Marathi for a cupboard.
\end{quote}

\subsection*{You'll want to know: \mr{जमीन}}
\textbf{\mr{जमीन}} — \emph{jamīn} — \textquotedblleft{}the ground, the floor\textquotedblright{}.

\begin{cousinweb}
\textbf{\mr{जमीन}} (\emph{jamīn}) is a loan, usually traced to Persian \emph{zamīn}, \textquotedblleft{}earth, ground\textquotedblright{}. Marathi writes the \emph{z} as \textbf{\mr{ज}}.
\end{cousinweb}

\begin{grammarlens}[title={What you've built}]
One more word for this chapter.
\end{grammarlens}

\subsection*{Your turn}
\begin{itemize}
\raggedright
  \item \emph{Say it:} \emph{jamīn}
  \item \emph{Say it:} \emph{jamīn}, once more
  \item \emph{Say it:} say \emph{jamīn}
  \item \emph{Recall:} say the Marathi for a bed, then the Marathi for a cupboard, then say \emph{jamīn} again
\end{itemize}

\begin{tcolorbox}[breakable,colback=teal!4,colframe=teal!35!black,title={Before you move on}]
What does \mr{जमीन} mean? (\textquotedblleft{}The ground, the floor\textquotedblright{}.) Say it once more.
\end{tcolorbox}
